
\textbf{Remarque 5.3.} Soient \(k\) un corps alg\'ebriquement clos, et \(X\) une surface propre et lisse sur \(k\). Alors le deuxi\`eme nombre de Betti \(b^2(X)\) peut \^etre interpr\'et\'e comme \'etant \(c_2(X)-2+2b^1(X)\), o\`u \(b^1(X)=2\dim(\Pic_{X/k})\), et \(c_2(X)\) = la deuxi\`eme classe de Chern = la caract\'eristique d'Euler-Poincar\'e = le nombre d'intersection de la diagonale de \(X\times X\) avec elle-m\^eme\,; ainsi \(b^2(X)\) a une interpr\'etation ind\'ependante de la cohomologie. L'in\'egalit\'e \(\rang(C^1_{\rm num}(X))\le b^2(X)\) dans cette forme pour \(\operatorname{car}(k)>0\) est due (avant la cohomologie \'etale) \`a IGUSA, qui a utilis\'e des r\'esultats profonds sur la structure du groupe fondamental mod\'er\'e d'une courbe alg\'ebrique moins quelques points.

\textbf{Th\'eor\`eme 5.4.} \emph{Soient \(k\) un corps alg\'ebriquement clos, et \(X\) un \(k\)-sch\'ema propre. Alors il existe un nombre fini de courbes ferm\'ees int\`egres \(Y_1,\ldots,Y_p\) dans \(X\), ayant la propri\'et\'e suivante~: pour qu'une partie \(\Lambda\) du sch\'ema \(\Pic_{X/k}\) soit quasi-compacte, il faut et il suffit que les nombres d'intersection \(\langle\lambda.Y_i\rangle\) des \(\lambda\in\Lambda\) restent born\'es. De plus, on peut prendre pour \(p\) le rang du groupe de N\'eron-Severi \(\Pic_{X/k}(k)/\Pic^0_{X/k}(k)\).}

En effet, comme \(\Pic^\tau_{X/k}\) est un sous-groupe ouvert quasi-compact de \(\Pic_{X/k}\) (4.7), pour que \(\Lambda\) soit quasi-compacte il faut et il suffit que l'ensemble correspondant de points de l'espace vectoriel \(\bigl[\Pic_{X/k}(k)/\Pic^\tau_{X/k}(k)\bigr]\otimes\Q\) soit fini. Mais la \(\tau\)-\'equivalence est \'egale \`a l'\'equivalence num\'erique (4.6 (i) \(\Longleftrightarrow\) (iii)), donc on peut choisir parmi les courbes dans \(X\) des courbes qui s'en vont sur une base de l'espace dual\,; d'o\`u l'assertion.

